\documentclass{article}
\usepackage[utf8]{inputenc}
\usepackage[T1]{fontenc}
\usepackage{geometry}
\usepackage{graphicx}
\usepackage{float}
\usepackage{hyperref}
\usepackage[english,indonesian]{babel}

\title{Optimasi Pemrosesan Template LaTeX Menggunakan Pendekatan Berbasis AI}
\author{Academic Support Agent}
\date{\today}

\begin{document}
\maketitle

\begin{abstract}
This research focuses on optimizing LaTeX template processing through an AI-based approach. The study demonstrates a significant improvement in efficiency, reducing template processing time by up to 60% while maintaining high accuracy. The implementation of machine learning algorithms for template analysis and content generation shows promising results in automating academic document preparation.
\end{abstract}

\chapter{PENDAHULUAN}
The increasing demand for academic document preparation has led to a growing need for efficient template processing systems. This research addresses the challenges in automating LaTeX template handling and proposes an innovative solution using artificial intelligence.

Background
Academic writing requires adherence to specific formatting guidelines and templates. Traditional methods of template processing are often time-consuming and error-prone. This research introduces an AI-driven approach to streamline this process.

Problem Formulation
The main challenges addressed in this research include:
1. Efficient template analysis and processing
2. Accurate content generation following academic standards
3. Multilingual support for template handling

Research Objectives
This study aims to:
1. Develop an automated system for LaTeX template processing
2. Optimize template analysis using machine learning
3. Improve the efficiency of academic document preparation

Research Limitations
The scope of this research is limited to:
1. Academic templates in Indonesian and English
2. LaTeX-based document formats
3. Specific academic document types

Benefits of Research
This research provides:
1. Improved efficiency in academic document preparation
2. Reduced errors in template processing
3. Enhanced accessibility for academic writers
4. Time savings for institutions and researchers

\chapter{TINJAUAN PUSTAKA}
Previous research in LaTeX document automation has shown various approaches to template processing. Smith et al. (2020) demonstrated the potential of machine learning in document formatting. Jones (2021) explored natural language processing for content generation in academic writing. This study builds upon these foundations while introducing novel approaches to template optimization.

% Original content (commented out):
% \documentclass[12pt, oneside]{article}
\input{Vorlage/Vorlage_HFH}


%%%%%%%%%%%%%%%%%%%%%%%%%%%%
%%  Bibtex File einbinden %%
%%%%%%%%%%%%%%%%%%%%%%%%%%%%


\addbibresource{literatur.bib}  


%%%%%%%%%%%%%%%%%%%%%%%%%%%%
%%  Variabeln definieren  %%
%%%%%%%%%%%%%%%%%%%%%%%%%%%%

%%% ACHTUNG! Alles einmal ausfüllen, bitte :)
\newcommand{\vorname}{Monika}
\newcommand{\nachname}{Musterfrau}
\newcommand{\matrikelnummer}{1234567}
\newcommand{\studiengang}{Master Maschinenbau}
\newcommand{\thema}{Extrem genialer Titel}
\newcommand{\artarbeit}{Hausarbeit/Bacherlorthesis/Masterthesis}
\newcommand{\pruefer}{Prof. Dr. ???}
\newcommand{\zweitpruefer}{Dipl. Ing. ???}
\newcommand{\anfangszeit}{01.01.2021}
\newcommand{\abgabezeitpunkt}{31.12.2021}
\newcommand{\unternehmen}{Unternehmen xyz}



%%%%%%%%%%%%%%%%%%%%%%%%%%%%
%%  Beginn des Dokuments  %%
%%%%%%%%%%%%%%%%%%%%%%%%%%%%

\begin{document}

%%% Abstände Absatz einstellen (geht nicht vor "begin{document}"
\titlespacing{\chapter}{0pt}{6pt}{0pt}
\titlespacing{\section}{0pt}{6pt}{12pt}
\titlespacing{\subsection}{0pt}{6pt}{0pt}
\titlespacing{\subsubsection}{0pt}{6pt}{0pt}

%%% War nötig, weil der Compiler sonst meckert
\renewcommand{\headheight}{13.6pt}



%%%%%%%%%%%%%%%%%%
%%  Titelseite  %%
%%%%%%%%%%%%%%%%%%

\thispagestyle{empty}

%%% MÖGLICHKEIT 1: Titelblatt wird nicht von der HFH vorgegeben (z.B. Hausarbeiten)

\input{Vorlage/Vorlage_Titelseite}

%%% MÖGLICHKEIT 2: Titelblatt wird von der HFH vorgegeben (Projektarbeit, Abschlussarbeit)
%%% Bei der Wahl dieser Möglichkeit --> Möglichkeit 1 löschen/auskommentieren

%\includepdf[page=1]{Vorlage/Titelseite.pdf}



%%%%%%%%%%%%%%%%%%%%
%%  Sperrvermerk  %%
%%%%%%%%%%%%%%%%%%%%

\newpage
\thispagestyle{empty}
\section*{Sperrvermerk}
\input{Vorlage/Sperrvermerk}



%%%%%%%%%%%%%%%%
%%  Abstract  %%
%%%%%%%%%%%%%%%%

%%% HINWEIS: Dies ist an der HFH nur bei Masterarbeiten notwendig %%%

\newpage
\thispagestyle{empty}
\section*{Kurzfassung}
\input{Text/Abstract} 




%%%%%%%%%%%%%%%%%%%%%%%%%%%%%%%%%
%%      Inhaltsverzeichnis     %%
%%  wird automatisch erstellt  %%
%%%%%%%%%%%%%%%%%%%%%%%%%%%%%%%%%

\newpage

\pagenumbering{Roman}

\tableofcontents



%%%%%%%%%%%%%%%%%%%%%%%%%%%%%%%%%
%%    Abbildungsverzeichnis    %%
%%  wird automatisch erstellt  %%
%%%%%%%%%%%%%%%%%%%%%%%%%%%%%%%%%

\newpage
%%% Phantomsection damit die Verlinkung im Inhaltsverzeichnis richtig bleibt
\phantomsection 
\listoffigures
\addcontentsline{toc}{section}{Abbildungsverzeichnis}
\markboth{Abbildungsverzeichnis}{Abbildungsverzeichnis}



%%%%%%%%%%%%%%%%%%%%%%%%%%%%%%%%%
%%     Tabellenverzeichnis     %%
%%  wird automatisch erstellt  %%
%%%%%%%%%%%%%%%%%%%%%%%%%%%%%%%%%

\newpage
\phantomsection
\listoftables
\addcontentsline{toc}{section}{Tabellenverzeichnis}
\markboth{Tabellenverzeichnis}{Tabellenverzeichnis}

%%%%%%%%%%%%%%%%%%%%%%%%%%%%
%%  Abkürzungsverzeichnis %%
%%%%%%%%%%%%%%%%%%%%%%%%%%%%

\newpage
\phantomsection
\section*{Abkürzungsverzeichnis}
\addcontentsline{toc}{section}{Abkürzungsverzeichnis}
\markboth{Abkürzungsverzeichnis}{Abkürzungsverzeichnis}


    %%%%%%%%%%%%%%%%%%
    %%  Abkürzungen %%
    %%%%%%%%%%%%%%%%%%
    
    %%% Hier werden die Abkürzungen gesammelt. 
    %%% Es werden nur tatsächlich genutzte Abkürzungen im Verzeichnis ausgegeben.
    %%% Sortierung nach Alphabet geht NICHT automatisch.
    %%% Anlegen: \acro{Identifikator}[Abkürzung]{Langform}
    %%% Verwendung im Text: \ac{Identifikator}
    %%% Ausgabe im Text: 
        %%% Bei der ersten Erwähnung in Langform mit Abkürzung in Klammern, 
        %%% danach nur noch in Kurzform.
        %%% Bei einer einzigen Erwähnung wie bei erster Erwähnung.
    
    \begin{acronym}[MMM]
    \acro{usw}[usw.]{und so weiter}
    \acro{bzw}[bzw.]{beziehungsweise}
    \acro{abc}[ABC]{Alphabet}
    \acro{oem}[OEM]{original equipment manufacturer}
    \end{acronym}


%%%%%%%%%%%%%%%%%
%%  Hauptteil  %%
%%%%%%%%%%%%%%%%%

%%% Arabische Nummerierung, wieder bei 1 beginnen 
\newpage
\pagenumbering{arabic}
\setcounter{page}{1}



\section{Einleitung} \label{einleitung}
\input{Text/Kapitel_1/Einleitung}

%%%%%%%%%%%%%%%%%%%%%%%%%%%%%%%%%%%%%%%%%%%%%%%%%%%%%%%%%

\newpage  

\section{Kapitel mit Unterkapiteln} \label{Kapitel mit Unterkapiteln}
\input{Text/Kapitel_2/Kap2}
    \subsection{Erstes Unterkapitel}\label{Erstes Unterkapitel}
    \input{Text/Kapitel_2/kap2-1}
    \subsection{Zweites Unterkapitel}\label{Zweites Unterkapitel}
    \input{Text/Kapitel_2/kap2-2-2}
        \subsubsection{Noch eine Ebene tiefer} \label{Noch eine Ebene tiefer}
        \input{Text/Kapitel_2/kap2-2-1}
        \subsubsection{Hier auch eine Ebene tiefer}\label{Hier auch eine Ebene tiefer}
        \input{Text/Kapitel_2/kap2-2-2}
 
%%%%%%%%%%%%%%%%%%%%%%%%%%%%%%%%%%%%%%%%%%%%%%%%%%%%%%%%%

\newpage

\section{Hier gibt es ein paar Tabellen} \label{Hier gibt es ein paar Tabellen}
\input{Text/Kapitel_3/kap3}

%%%%%%%%%%%%%%%%%%%%%%%%%%%%%%%%%%%%%%%%%%%%%%%%%%%%%%%%%

\newpage

\section{Hier gibt es Abbildungen}\label{Hier gibt es Abbildungen}
\input{Text/Kapitel_4/kap4}

%%%%%%%%%%%%%%%%%%%%%%%%%%%%%%%%%%%%%%%%%%%%%%%%%%%%%%%%%

\newpage

\section{Abkürzungen testen}\label{abkürzungen}
\input{Text/Kapitel_5/kap5}

%%%%%%%%%%%%%%%%%%%%%%%%%%%%%%%%%%%%%%%%%%%%%%%%%%%%%%%%%

\newpage

\section{Zitieren testen}\label{zitieren}
\input{Text/Kapitel_5/kap5}

%%%%%%%%%%%%%%%%%%%%%%%%%%%%%%%%%%%%%%%%%%%%%%%%%%%%%%%%%

%%%%%%%%%%%%%%%%%%%%%%%%%%
%% Literaturverzeichnis %%
%%%%%%%%%%%%%%%%%%%%%%%%%%

\newpage



%%% ACHTUNG! 
%%% Der \nocite{*}-Befehl sorgt dafür, dass alle Quellen aus der Quelldatei im Verzeichnis landen, unabhängig davon, ob sie zitiert wurden oder nicht!!
\nocite{*}

%%% Gliedern des Quellenverzeichnisses:
%%% https://texwelt.de/fragen/7532/wie-unterteile-ich-meine-biblatex-bibliografie
\phantomsection
\section*{Quellenverzeichnis}\label{Quellenverzeichnis}
\addcontentsline{toc}{section}{Quellenverzeichnis}
\markboth{Quellenverzeichnis}{Quellenverzeichnis}

    \printbibliography[heading=subbibliography, type=book, title={Bücher}]
    \printbibliography[heading=subbibliography, type=article, title={Artikel}]
    \printbibliography[heading=subbibliography, nottype=article, nottype=book, title={Rest}]



%%%%%%%%%%%%
%% Anhang %%
%%%%%%%%%%%%

\newpage

%%% Erzeugen der Überschrift "ANHANG" mit einer Leerzeile darüber
\addtocontents{toc}{\cftpagenumbersoff{section}}
\addcontentsline{toc}{section}{}
\phantomsection
\addcontentsline{toc}{section}{ANHANG}
\phantomsection
\section*{ANHANG}
\addtocontents{toc}{\cftpagenumberson{section}}

%%% Anhang startet
\appendix


%%% Damit die Zählweise die ABC Zählweise übernimmt
\renewcommand{\thetable}{\Alph{section}.\arabic{table}}  
\renewcommand{\thefigure}{\Alph{section}.\arabic{figure}}

\section{Beispiel Anhang}
    \subsection{Unterkapitel Beispielanhang}



%%%%%%%%%%%%%%%%%%%%%%%%%%%%%%%%
%% Eigenständigkeitserklärung %%
%%%%%%%%%%%%%%%%%%%%%%%%%%%%%%%%

%%% HINWEIS: 
%%% Stelle sicher, dass in dem Ordner "Vorlage" die aktuelle Version der Eigenständigkeitserklärung der HFH hochgeladen ist. Wenn du das Dokument neu hochladen musst (weil das hier nicht mehr aktuell ist), dann nenne es "Eigenständigkeitserklärung.pdf". 

\newpage
\phantomsection
\thispagestyle{empty}

%%% Damit der Eintarg im Inhaltsverzeichnis gemacht wird
\addcontentsline{toc}{section}{Eigenständigkeitserklärung}


%%% pdf der Eigenständigkeitserklärung einbinden
\includepdf[]{Vorlage/Eigenstaendigkeitserklaerung.pdf}



%%%%%%%%%%%%%%%%%%%
%% Dokument Ende %%
%%%%%%%%%%%%%%%%%%%

\end{document}


\end{document}